\section{Sources, conservation et reproduction}
Le supplément numérique conserve le document cumulatif intégral, ses trente-neuf sources, les programmes de vérification et les reçus correspondants. La narration approuvée et sa continuation sont maintenues littéralement dans cette archive. La présente édition transforme leur énonciation conversationnelle en exposition académique et consigne la correspondance entre sections, sans remplacer les originaux.

Les antécédents proviennent des sources de la série du 21 septembre 2026. Le noyau formel est incorporé à partir de sa clôture documentaire dans l'article constitutif ; les développements de l'action, de l'incidence, de la masse et de la température sont conservés avec leurs démonstrations. Le manifeste indique chaque fichier, intervalle extrait et empreinte. Ces données documentent la provenance et la conservation, tandis que les preuves du corps du texte fixent la portée mathématique des résultats.

Les programmes de composition des constantes, d'hybridation et de compatibilité confrontent des identités exactes et des exemples finis. Les diagnostics numériques sont distingués des contrôles rationnels. Les preuves de monotonie, de convergence, d'inversion et de passage à la limite restent développées dans le texte : les exemples exécutables complètent ces démonstrations.

La compilation utilise LuaLaTeX. Les sources incluses et leurs figures permettent de reproduire le PDF sans dépendre des répertoires originaux de l'auteur. Le fichier d'instructions du paquet spécifie l'ordre de compilation et le mode d'exécution des vérifications.
